\documentclass[11pt]{amsart}
\usepackage[T1]{fontenc}
\usepackage{lmodern}
\usepackage{amsmath,amssymb,amsthm}
\usepackage[margin=1.1in]{geometry}
\usepackage{booktabs}
\usepackage{enumitem}
\usepackage{hyperref}

\newtheorem{theorem}{Theorem}[section]
\newtheorem{lemma}[theorem]{Lemma}
\newtheorem{proposition}[theorem]{Proposition}
\newtheorem{corollary}[theorem]{Corollary}
\newtheorem{conjecture}[theorem]{Conjecture}
\theoremstyle{definition}
\newtheorem{definition}[theorem]{Definition}
\newtheorem{remark}[theorem]{Remark}
\newtheorem{example}[theorem]{Example}

\newcommand{\cD}{\mathcal D}
\newcommand{\lean}[1]{\quad{\footnotesize[Lean: \texttt{#1}]}}
\newcommand{\MT}{\mathrm{MT}}
\emergencystretch=3em

\title[Multi-token steering]{Multi-token steering:\\ the triangle base from $k$-token schedules}
\author{Gallai project (cloud session \texttt{claude/multi-token})}
\date{28 November 2026 (working draft)}

\begin{document}
\begin{abstract}
The triangle base~\cite{DT} asks for a Lov\'asz decomposition of an auxiliary graph $M$ with one
disjunctive end-path condition per pair of a matching; it implies Test1, one of the two tests to
which Gallai's conjecture is reduced in Lean. The frontier is $d(z)=6$: two pairs, and in the
dense T(6)-i residual no proof is known. Two-token steering~\cite{TRN} settles one pair by hanging a
leaf and moving two tokens. We extend it to $k$ pairs at once. \emph{Status labels: the definitions, lemmas and theorems of
Sections~\ref{sec:sched}--\ref{sec:tb} and Proposition~\ref{prop:nminus2} are proved and formalised in Lean
(Remark~\ref{rem:general} is not formalised, and Example~\ref{ex:short} is a computed example);
Theorem~\ref{thm:cover}, Proposition~\ref{prop:potential}, Proposition~\ref{prop:aoddlift} and the instance in
Proposition~\ref{prop:aoddobs} are proved on paper; the counts of Sections~\ref{sec:cover}--\ref{sec:aodd}
are computations with stated scope.}
\begin{itemize}
\item \emph{$k$-token schedules and the transfer theorem} (Lean). Hang a leaf $\ell_i$ at one vertex
$y_i$ of every pair; the tokens are the even vertices of the current graph. A schedule uses three
moves: a Fan step of a token onto an odd vertex outside $\{\ell_i,x_i\}$ for the live tails that,
after the deletion of the edge, is adjacent to no other token; a \emph{drop} of a leaf whose stem has become odd; and the \emph{suppression} of a token of
degree two. It ends with at most one token and one live tail, where the \emph{multi-tail leaf lemma}
applies. A schedule gives a Lov\'asz decomposition of $M$ with $x_i\notin Q(y_i)$ for every $i$
(Theorems~\ref{thm:transfer} and~\ref{thm:tb}), hence a triangle base and Test1
(Corollary~\ref{cor:test1}).
\item \emph{Coverage} (theorem on paper; the count is computed). An explicit certificate family
(Theorem~\ref{thm:cover}) covers $882$ of the $10{,}118$ roots with data at order~$9$. No structural
condition covering a whole dense residual shape exists for this move set: every shape of both residual
families has roots of degree $n-2$, which have no schedule (Proposition~\ref{prop:nminus2}).
\item \emph{Computation.} At order~$9$ ($d(z)=6$, all $1{,}819$ $3$-connected $J$) some data choice
and orientation has a schedule at $7{,}529$ of the $10{,}118$ roots ($74\%$); all $76{,}474$
schedules found were executed and checked by an independent program, including the lift to Test1. On
the residual families of~\cite{T6}, at node budget $2\cdot10^6$ (and $2\cdot10^7$ for the roots left
open), the proportion is $72$ of $180$ roots ($|O|=8$) and $111$ of $144$ ($|O|=10$, with a $C_4$
subfamily at $|O|=8$); on the eight E28 hard graphs it is $10$ of $48$ roots.
\item \emph{The residual} (paper and computation). A schedule needs $2k-1$ \emph{removals}
(retirements, suppressions, and drops that isolate the stem), each at a vertex of degree at most two
(Proposition~\ref{prop:potential}). If $r$ has fewer than $k$
non-neighbours, at most $k-1$ removals are possible, so there is no schedule
(Proposition~\ref{prop:nminus2}, Lean): at $d(z)=6$ this excludes every root of degree $n-2$, which is
$1{,}597$ of the $2{,}589$ unscheduled roots at order~$9$, and roots of every residual shape. In the dense residual, $42$ of
the $53$ other roots that the search proves unscheduled are \emph{frozen}: no reachable state allows a
single removal.
\item \emph{The all-odd lift} (Section~\ref{sec:aodd}). On the all-odd base of~\cite[\S31]{SYN}, where
Test1 needs a Lov\'asz decomposition of an all-odd graph with one fixed and $k-1$ disjunctive
conditions, steering does worse: $2{,}101$ of $8{,}340$ roots at order~$9$ ($25\%$). The obstruction
survives: no root with at most $k-2$ non-neighbours has a schedule (Proposition~\ref{prop:aoddobs}).
\end{itemize}
\end{abstract}
\maketitle

\section{Introduction}\label{sec:intro}

\subsection*{The triangle base} Let $J$ have exactly one even vertex $z$, with $d(z)=2m$. Triangle
data at a root $r\in N(z)$ are a vertex $w\in N(z)-r$ and a perfect matching
$\{a_1b_1,\dots,a_kb_k\}$, $k=m-1$, of $J[N(z)-\{r,w\}]$ by edges of $J$. Put
\[M=J-zw-\{za_i,zb_i,a_ib_i:i\le k\}.\]
In $M$ the vertex $z$ is a leaf at $r$, $w$ is the only even vertex, and $a_i\not\sim b_i$. A
\emph{triangle base} is a Lov\'asz decomposition $\cD$ of $M$ with, for every $i$,
\begin{equation}\label{eq:disj}
a_i\notin V(Q(b_i))\quad\text{or}\quad b_i\notin V(Q(a_i)),
\end{equation}
where $Q(v)$ is the path ending at $v$. A triangle base gives Test1 at $(J,z,r)$~\cite[Thm.~2.1]{DT}
(Lean \texttt{test1At\_of\_triangleBase}). For $d(z)=4$ there is one pair, and Test1 at $d(z)=4$ was
already known (two-step descent and \texttt{emc\_four}~\cite[\S28]{SYN}); the frontier is $d(z)=6$, two pairs, in particular the dense T(6)-i residual, where
$J[N(z)]=K_6-D$ with $D\in\{\emptyset,K_2,P_3,K_3,C_4\}$~\cite{T6}. There the triangle base holds
in all data~\cite[\S6]{DT}, without proof.

\subsection*{Steering} Neighbour avoidance~\cite{NA} turned ``$Q(y)$ avoids $x$'' into a property
of vertex sets: hang a leaf $\ell$ at $y$; then the property is ``no path contains both $\ell$ and
$x$'', and Fan's lemma with vertex sets keeps it through every step whose target is not $\ell$ or
$x$. Two-token steering~\cite[\S6]{TRN} runs Lov\'asz's induction with two even vertices, $g$ and
$y$, choosing the targets in advance so that Fan's inequality never becomes tight.

\subsection*{This paper} For $k$ pairs we orient each pair, $y_i\in\{a_i,b_i\}$ and $x_i$ the
other vertex, hang a leaf $\ell_i$ at every $y_i$, and play with $k+1$ tokens. Three points are new.
\begin{enumerate}
\item The tokens are the even vertices of the current graph. A Fan step needs its target to be
non-adjacent to \emph{every} other token (Lemma~\ref{lem:fan}).
\item Two further moves are sound: a live tail whose stem has become odd can be \emph{dropped}, by
deleting its leaf edge, which comes back as a one-edge path (Lemma~\ref{lem:steps}(b)); and a token
of degree two with non-adjacent neighbours can be suppressed (Lemma~\ref{lem:steps}(c)). Drops are
what lets us end with the proved multi-tail leaf lemma: at order~$9$, schedules with drops ending
at one live tail reach $7{,}529$ roots, schedules without drops ending at any number of live tails
(which would need Conjecture~\ref{conj:mt}) reach $7{,}558$.
\item The end is the \emph{multi-tail leaf lemma} (Lemma~\ref{lem:mtl}): with at most one token and
at most one live tail, Lov\'asz's theorem or the leaf-tail lemma of~\cite{TRN} applies. The
general multi-tail statement, with any number of live tails, is Conjecture~\ref{conj:mt}; it holds
in all our tests but adds little coverage ($91$ roots at order~$9$).
\end{enumerate}

\section{Tails and the moves}\label{sec:sched}

All graphs are finite and simple. For a path decomposition $\cD$ let $b_\cD(v)$ be the number of
path ends at $v$. A \emph{parity decomposition} of $H$ has $b(v)=d_H(v)\bmod2$ for every $v$; when
$H$ has at most one vertex of positive even degree it is a Lov\'asz decomposition, and one exists
(Lov\'asz~\cite{L68}; Lean \texttt{lovasz\_ends}). A vertex is \emph{even} if its degree is even and
positive, \emph{odd} if its degree is odd.

\begin{definition}[tails]\label{def:tails}
A \emph{tail} is a triple $t=(c,s,x)$ of vertices: a \emph{leaf} $c$, its \emph{stem} $s$, and the
vertex $x$ to be avoided. It is \emph{live} in $K$ if $cs\in E(K)$. Let $T$ be a finite set of
tails.
\begin{itemize}
\item[(TS)] \emph{Static conditions}: $x\notin\{c,s\}$ for every tail, and for distinct tails
$t,t'$ the leaf $c'$ and the avoided vertex $x'$ of $t'$ lie outside $\{c,s\}$.\lean{TailSys}
\item[(TI)] \emph{Invariant} in $K$: for every live tail, $N_K(c)=\{s\}$, $x$ is odd and
$x\not\sim s$.\lean{TailInv}
\item[(TG)] $K$ is \emph{good} if it has a parity decomposition in which, for every live tail, no
path contains both $c$ and $x$.\lean{TailGood}
\end{itemize}
Let $X(K)=\bigcup\{\{c,x\}:(c,s,x)\in T\text{ live in }K\}$.
\end{definition}

Since $c$ is a leaf, the path through $c$ ends at $c$, so (TG) says that for every live tail the
path ending at $c$ avoids $x$. It is a property of vertex sets.

\begin{definition}[$k$-token schedule]\label{def:sched}
The \emph{tokens} of $K$ are its even vertices, $\tau(K)$. A \emph{move} replaces $K$ by $K'$:
\begin{enumerate}[label=(\Alph*)]
\item[(M)] \emph{token move}: a token $p$ and an edge $pu$ with $u$ odd, $u\notin X(K)$, and no
neighbour of $u$ in $K-pu$ a token other than $p$; then $K'=K-pu$. The token goes to $u$; it
\emph{retires} if $u$ is isolated in $K'$.
\item[(D)] \emph{drop}: a live tail $(c,s,x)$ whose stem $s$ is odd; then $K'=K-cs$. The stem becomes
a token unless it is isolated in $K'$.
\item[(S)] \emph{suppression}: a token $p$ with $N_K(p)=\{a,b\}$, $a\not\sim b$, neither $a$ nor
$b$ the leaf of a tail, and $\{a,b\}\neq\{x,s\}$ for every live tail; then $K'=K-pa-pb+ab$.
\end{enumerate}
A \emph{schedule} for $(K,T)$ is a sequence of moves after which at most one token and at most one
live tail are left.\lean{SteerK}
\end{definition}

In Lean the schedule is the inductive predicate \texttt{SteerK T K $\tau$}, with the token set
$\tau$ carried explicitly; the transfer theorem assumes $\tau=\tau(K)$ at the start, and each
constructor states the new token set: $\tau-p$, plus $u$ unless $u$ is isolated, after (M); $\tau$,
plus $s$ unless $s$ is isolated, after (D); $\tau-p$ after (S) (Lemma~\ref{lem:inv}).

\begin{lemma}[Fan's step with many even vertices]\label{lem:fan}
Let $p$ be even in $K$, $u$ an odd neighbour of $p$, and suppose that in $K-up$ the vertex $u$ has
no even neighbour of $K$ other than $p$. Every parity decomposition $\cD'$ of $K-up$ gives a parity
decomposition $\cD$ of $K$, every path of which lies in $V(P)\cup\{u\}$ for some $P\in\cD'$.
\lean{fan\_step\_multi}
\end{lemma}

\begin{proof}
In $K'=K-up$ the vertex $p$ is odd, so $b_{\cD'}(p)=1$, and $u$ is even or isolated, so
$b_{\cD'}(u)=0$. A neighbour $v$ of $u$ in $K'$ with $b_{\cD'}(v)=0$ would be even in $K'$, hence
$v\neq p$ and $v$ even in $K$ (its degree is unchanged), which is excluded. So $Z_{\cD'}(u)=\emptyset$
and $b_{\cD'}(p)=1>|Z_{\cD'}(u)|$. Fan's lemma with vertex sets~\cite{Fan05},~\cite[Lemma~2.2]{RD}
(Lean \texttt{fan\_single\_supp}) gives $\cD$ with $b_\cD=b_{\cD'}-\mathbf 1_p+\mathbf 1_u$ and the
covering property. This is $0$ at $p$, $1$ at $u$, and $d\bmod2$ elsewhere.
\end{proof}

\begin{lemma}[bookkeeping]\label{lem:inv}
Let $T$ satisfy (TS) and $K$ satisfy (TI).
\begin{enumerate}[label=(\alph*)]
\item After (M), (D) or (S) the token set is as stated above, and (TI) holds in
$K'$.\lean{evenSet\_move, evenSet\_drop, evenSet\_supp, tailInv\_move, tailInv\_drop, tailInv\_supp}
\item Every live tail of $K$ stays live in $K'$, except the dropped one after (D).\lean{live\_move,
live\_drop, live\_supp}
\end{enumerate}
\end{lemma}

\begin{proof}
(a) Tokens. In (M), $p$ loses one edge and becomes odd; $u$ is odd and loses one edge, so it is even
unless isolated; other degrees are unchanged. In (D), $c$ becomes isolated and the odd stem $s$
becomes even unless isolated. In (S), $p$ becomes isolated and the degrees of $a,b$ are unchanged.
(TI): edges are only deleted, except $ab$ in (S). The leaf $c$ of a live tail is neither $p$ (a
token; $d(c)=1$) nor $u$ ($u\notin X$), nor $a,b$ (excluded in (S)), nor the leaf or stem of the
dropped tail (by (TS)); so $N(c)$ is unchanged. The vertex $x$ is odd, hence not a token, and $x\ne
u$; in (D) $x$ lies outside $\{c_0,s_0\}$ by (TS); in (S) the degrees of $a,b$ are unchanged. So $x$
stays odd. Finally $x\not\sim s$ survives deletions, and the new edge $ab$ of (S) is not $xs$.
(b) The edge $cs$ of a live tail is not $pu$, $pa$ or $pb$ (as $c\notin\{p,u,a,b\}$), and in (D) it
is not the dropped edge $c_0s_0$ unless the tail is the dropped one.
\end{proof}

\begin{lemma}[the three steps]\label{lem:steps}
Let $T$ satisfy (TS), $K$ satisfy (TI), and $K'$ arise from $K$ by one move. If $K'$ is good, so
is $K$.
\begin{enumerate}[label=(\alph*)]
\item (M): by Lemma~\ref{lem:fan}; the new paths lie in old paths plus $u\notin X(K)$.
\lean{tailGood\_move}
\item (D): add the one-edge path $cs$ to a good decomposition of $K-cs$. The stem $s$ was even or
isolated in $K-cs$ and is odd in $K$; the path $cs$ misses $x$ and the vertices of every other tail
(by (TS)); no other path contains the isolated vertex $c$.\lean{tailGood\_drop}
\item (S): subdivide the edge $ab$, on its path of a good decomposition of $K-pa-pb+ab$, by $p$.
The ends do not change, $p$ gets none, and the new path lies in the old one plus $p\notin X$.
\lean{tailGood\_supp, splitOff\_supp}
\end{enumerate}
\end{lemma}

\begin{lemma}[the multi-tail leaf lemma]\label{lem:mtl}
Let $T$ satisfy (TS), $K$ satisfy (TI), and suppose $K$ has at most one even vertex and at most one
live tail. Then $K$ is good.\lean{tailGood\_fin}
\end{lemma}

\begin{proof}
With no live tail this is Lov\'asz's theorem. With one live tail $(c,s,x)$, $c$ is a leaf at $s$,
$x\notin\{c,s\}$ is odd and $x\not\sim s$, so the leaf-tail lemma~\cite[Lemma~6.1]{TRN} (Lean
\texttt{avoid\_leafTail}; the stem may be the even vertex) gives a Lov\'asz decomposition whose
path ending at $c$ avoids $x$. Since $c$ is a leaf, every path through $c$ ends at $c$, so no path
contains both $c$ and $x$.
\end{proof}

\begin{theorem}[the transfer theorem]\label{thm:transfer}
Let $T$ satisfy (TS) and $K$ satisfy (TI). If $(K,T)$ has a $k$-token schedule, then $K$ has a
parity decomposition in which, for every tail live in $K$, no path contains both its leaf and its
avoided vertex.\lean{steerK\_good}
\end{theorem}

\begin{proof}
Induction on the schedule. At the end Lemma~\ref{lem:mtl} applies. Each move keeps (TI) and gives
the right token set (Lemma~\ref{lem:inv}), so the induction hypothesis applies to $K'$, and
Lemma~\ref{lem:steps} carries goodness back to $K$.
\end{proof}

\begin{remark}\label{rem:general}
Nothing in Theorem~\ref{thm:transfer} refers to the triangle base: $K$ is arbitrary, and the number
of tokens and of tails is arbitrary. With one tail and the two tokens $g,y$ of~\cite[\S6]{TRN},
schedules without (D) and (S) are the two-token schedules of~\cite[Def.~6.3]{TRN}: with two even
vertices a target other than the second token is odd, and the end there, a retirement, leaves one
token and the one tail.
\end{remark}

\section{The triangle base}\label{sec:tb}

Fix triangle data $(J,z,r,w,\{a_ib_i\})$ with $z$ the only even vertex of $J$, and an orientation:
for each $i$ a vertex $y_i\in\{a_i,b_i\}$, and $x_i$ the other. Let $M^+$ be $M$ with a new leaf
$\ell_i$ at each $y_i$, and $T=\{(\ell_i,y_i,x_i):i\le k\}$. The tokens of $M^+$ are $w$ and the
$y_i$ (when $d_J(w)\ge3$). In Lean $M^+$ is \texttt{leaves M Y} on $V\oplus V$, with $\ell_i$ the
copy \texttt{inr}~$y_i$, and the orientation is the order of the coordinates of each pair.
\lean{triLeaf, triTails}

\begin{theorem}[schedules give triangle bases]\label{thm:tb}
If $(M^+,T)$ has a $k$-token schedule, then $M$ has a Lov\'asz decomposition with $x_i\notin
V(Q(y_i))$ for every $i$; in particular $J$ has a triangle base at $r$.
\lean{triangleBase\_of\_steerK}
\end{theorem}

\begin{proof}
(TS) holds because the $2k$ pair vertices are distinct. (TI) holds in $M^+$: $\ell_i$ is a leaf
at $y_i$; $d_{M^+}(x_i)=d_M(x_i)$ is odd, since $x_i\ne w$ is not a leaf-bearer; and $x_iy_i\notin
E(M)$.\lean{triTails\_sys, triTails\_inv} By Theorem~\ref{thm:transfer}, $M^+$ has a parity
decomposition $\cD^+$ in which no path contains both $\ell_i$ and $x_i$. Delete the leaves (Lean
\texttt{Retract.delete\_supp} for \texttt{leavesRetract}). A vertex $v\notin\{y_i\}$ keeps its ends.
At $y_i$, which is even in $M^+$ and odd in $M$, $b_{\cD^+}(y_i)=0$ and the new count is at most
$1$ and odd, hence $1$. So the result $\cD$ is a Lov\'asz decomposition of $M$. Suppose the path $Q$
of $\cD$ ending at $y_i$ contains $x_i$. It is the trace of a path $P\in\cD^+$, which contains
$x_i$, hence not $\ell_i$. Then every edge of $P$ at $y_i$ is an edge of $M$, so $P$ has exactly one
edge at $y_i$ and ends there, against $b_{\cD^+}(y_i)=0$. So $x_i\notin V(Q(y_i))$, which is the
disjunct of~\eqref{eq:disj} with $y_i$ in the role of $b_i$ or $a_i$.
\end{proof}

\begin{corollary}[Test1]\label{cor:test1}
If $d_J(r)\ge3$ and some data choice at $r$ and some orientation have a schedule, Test1 holds at
$(J,z,r)$.\lean{test1At\_of\_steerK}
\end{corollary}

\begin{example}\label{ex:short}
$J=$ \texttt{HsaBEy\textasciitilde} (order $9$, $d(z)=6$), $z=7$, $r=3$, $w=1$, pairs oriented as
$(y,x)=(4,0),(8,5)$. Here $N_J(w)=\{z,0,6\}$ with $0\not\sim6$, and $d_J(4)=3$. The shortest schedule
(iterative deepening, \texttt{MAXLEN=8}), with four moves after the suppression, is
\[
\mathrm S\,1{:}\,0\text{-}6,\quad 8\to3,\quad 3\to z,\quad 4\to8,\quad \mathrm D\,(\ell_1,4,0).
\]
The token $w$ is suppressed; the token at $y_2=8$ moves to $r$ ($r\not\sim4$) and retires at the leaf
$z$; the token at $y_1=4$ moves onto $8$, which is odd again and has no token neighbour; then $4$ is
odd with only its leaf left, and the drop isolates it. One token ($8$) and one live tail
$(\ell_2,8,5)$ remain.
\end{example}

\section{Coverage}\label{sec:cover}

\begin{proposition}[the potential]\label{prop:potential}
Let $\Phi(K)=|\tau(K)|+\#\{\text{live tails}\}$. A token move that does not retire and a drop whose
stem is not isolated leave $\Phi$ unchanged; a retirement, a suppression and a drop that isolates
the stem lower it by one. In $M^+$, $\Phi=2k+1$ when $d_J(w)\ge3$, so a schedule contains at least $2k-1$ moves of the
second kind, each at a vertex of degree at most $2$ (a vertex of degree $1$ entered by a token, a
token of degree $2$, or a stem with only its leaf). If $\delta(J)\ge3$, the only vertex of degree $1$
in $M$ outside the pairs is $z$, and a token can retire there only from $r$.
\end{proposition}

\begin{proof}
Immediate from Lemma~\ref{lem:inv}(a).
\end{proof}

So at $d(z)=6$ three removals are needed, and in a dense $M$ low degrees arise only after tokens
have passed through a vertex (each pass deletes two of its edges). This is why schedules in the
dense residual are long, and why they fail when the root is hard to enter.

\begin{theorem}[a certificate family]\label{thm:cover}
Let $d(z)=6$, let $(r,w,\{x_1y_1,x_2y_2\})$ be oriented triangle data, and suppose:
\begin{enumerate}[label=(\Alph*)]
\item \emph{$w$ can be removed}: either (A1) $d_J(w)=3$, $N_J(w)=\{z,p,q\}$, $p\not\sim q$ and
$\{p,q\}$ is not a pair; or (A2) $wr\in E(J)$ and $r\not\sim y_1,y_2$.
\item \emph{$y_1$ has degree three and its token can be removed}: $N_J(y_1)=\{z,x_1,v\}$ with
$v\notin\{w,x_2,y_2\}$, $v\not\sim y_2$, and $\{p,q\}\ne\{v,y_2\}$ under (A1); and either (B1)
$v=r$ and (A1) holds, or (B2) $v\notin N[z]$, $d_J(v)=3$, and the two neighbours $s,t$ of $v$
other than $y_1$ in the graph obtained after step~1 of the proof below (that is, after the
suppression of (A1) if it is used) are non-adjacent there, with $\{s,t\}\ne\{x_2,y_2\}$.
\end{enumerate}
Then $(M^+,T)$ has a schedule; so $J$ has a triangle base at $r$, and if $d_J(r)\ge3$ (for example
if $J$ is $3$-connected) Test1 holds at $(J,z,r)$.
\end{theorem}

\begin{proof}
The tokens of $M^+$ are $w,y_1,y_2$; $X=\{\ell_1,\ell_2,x_1,x_2\}$. Write $K$ for the current graph.
\begin{enumerate}
\item Remove $w$. Under (A1), $d_M(w)=2$ with $N_M(w)=\{p,q\}$; neither is a leaf, $p\not\sim q$,
and $\{p,q\}$ is no pair $\{x_i,y_i\}$: suppress $w$. Under (A2), $r$ is odd, $r\notin X$, and after
deleting $wr$ no neighbour of $r$ is a token (the other tokens are $y_1,y_2$, and $r\not\sim
y_1,y_2$): move $w\to r$. Now $r$ is a token, $z$ is a leaf at $r$ and $z\notin X$: $r\to z$
retires.
\item Move $y_1\to v$. In $M$, $N_M(y_1)=\{v\}$. The vertex $v$ is odd in $K$ (it is not $w$, and
its degree is unchanged by step 1, also when $v\in\{p,q\}$), $v\notin X$, and after deleting $y_1v$
its only possible token neighbour is $y_2$, excluded by $v\not\sim y_2$ and, under (A1), by
$\{p,q\}\ne\{v,y_2\}$.
\item Remove the token at $v$. Under (B1), $v=r$ and $z$ is still a leaf at $r$: $r\to z$ retires.
Under (B2), $v$ has degree $2$ in $K$ with non-adjacent neighbours $s,t$, neither a leaf (since $v\ne
y_1,y_2$) and $\{s,t\}$ not $\{x_2,y_2\}$, nor $\{x_1,y_1\}$ (as $y_1\not\sim v$ now): suppress
$v$.
\item Drop $(\ell_1,y_1,x_1)$: $y_1$ has only its leaf, so it is odd, and the drop isolates it.
\end{enumerate}
One token, $y_2$, and one live tail, $(\ell_2,y_2,x_2)$, remain. Apply Theorem~\ref{thm:tb}, and
Corollary~\ref{cor:test1} when $d_J(r)\ge3$. The hypothesis is needed: under (A1) and (B2) the root
may be a leaf of $J$, and then Test1 fails although the schedule exists (audit A36, M2:
\texttt{Lsech?H?GA\_B?@}, $\kappa=1$, $r=1$ with $N(r)=\{z\}$).
\end{proof}

At order~$9$ the hypotheses hold, for some data choice, labelling and orientation, at $882$ of the
$10{,}118$ roots with data ($6$ of $40$ at order~$7$); \texttt{pattern.py} checks them from $J$ (at roots
with $d_J(r)\ge3$) and writes the certificate, and the independent executor of Section~\ref{sec:comp} verified all $882$.
The theorem is a sparse-graph statement: it needs vertices of degree three. In the dense residual
$\delta(J-z)\ge6$ and no such certificate exists.

\begin{remark}[no theorem for a whole residual shape]\label{rem:noshape}
For this move set no condition can cover a whole residual shape. In both residual families every
shape $D\in\{\emptyset,K_2,P_3,K_3,C_4\}$ has roots of degree $n-2$ ($15,10,5,9,3$ of them in
\texttt{resid\_O8}; $1,4,4,3,2$, and $4$ for $C_4$ with $\mathrm{TW}=1$, in \texttt{resid\_O10}), and these
have no schedule by Proposition~\ref{prop:nminus2} (Lean). In addition, at $|O|=8$ every shape has
roots of other degrees at which the exhaustive search (no budget hit for any data choice or
orientation) finds no schedule (Table~\ref{tab:resid}). These are computed. Our search memoises failed
states by a $128$-bit hash of the current graph; the exact search of audit A36, without hashing,
confirms every such verdict.
\end{remark}

\section{Computation}\label{sec:comp}

\subsection*{Tools} All scripts are in \texttt{manuscripts/multi-token/scripts}.
\begin{itemize}
\item \texttt{mtok.cpp}: depth-first search with memoised failures over the moves of
Definition~\ref{def:sched} (suppression with \texttt{SUPP=1}), for every root $r$ with $d(r)\ge3$,
every data choice $(w,M_0)$ and every orientation; per root the verdict is ``schedule'',
``none'' (every search finished without a schedule) or ``undecided'' (some search hit the node
budget). \texttt{MAXLEN} gives shortest schedules by iterative deepening.
\item \texttt{mt\_verify.py}, the independent executor. It shares no code with \texttt{mtok.cpp}.
For each schedule it rebuilds $J$, checks the data, rebuilds $M^+$, replays the moves re-deriving
their legality from Definition~\ref{def:sched}, checks the end condition, solves the bottom graph
with the oracle \texttt{route\_ext} (query \texttt{WS}) and re-checks that decomposition, undoes
the moves (a drop by the path $\ell_iy_i$, a suppression by subdividing $ab$, a token move by its
own implementation of Fan's forced chain, checking $b(p)=1$ and $Z(u)=\emptyset$), checks after every
step that it has a parity decomposition meeting every live $\{\ell_i,x_i\}$ at most once, deletes the
leaves, checks the Lov\'asz decomposition of $M$ and $x_i\notin Q(y_i)$, and finally lifts it to a
decomposition of $J-z$ with \texttt{tb.py}'s \texttt{lift} (direct-test1, Steps 1--2) and asserts
with \texttt{verify\_test1} that it is an edge partition into at most $(|J|-1)/2$ paths with at least
two ends at $r$. (Before commit \texttt{28376e6} the boolean returned by \texttt{verify\_test1}, which
carries the last two conditions, was not asserted; audit A36, M4. Every verification was rerun after
the fix and passes.)
\item A tamper test (\texttt{research/audit/A36-scripts/tamper.py}, mode \texttt{tamper}, seeds
$1,2,3$; output \texttt{data/runs/J7\_tamper\_seed*.txt}): each of the $130$ order-$7$ schedules is
corrupted once at random per seed (first or last move removed, a target changed, two adjacent moves
swapped, a wrong drop index, a move turned into a suppression), and the executor's verdict is compared
with an independent replayer of Definition~\ref{def:sched}. Of the $390$ corruptions the executor
accepted $38$, all of them valid schedules ($33$ swaps of commuting moves and $5$ valid suppressions),
and rejected the $352$ others, all invalid.
\item \texttt{pattern.py} (Theorem~\ref{thm:cover}), \texttt{firstmove.py} (legal first moves),
\texttt{filt.py} (the order-$7$ and order-$9$ inputs from \texttt{stl/scripts/one\_even\_*.g6}).
\end{itemize}

\begin{table}[ht]
\caption{Orders $7$ and $9$, $d(z)=6$, every $3$-connected $J$ with $z$ the only even vertex (no
node budget; all searches finished). ``Data'' counts data choices $(r,w,M_0)$; a data choice is
scheduled if some orientation is.}\label{tab:small}
\small
\begin{tabular}{@{}lrrrrr@{}}
\toprule
order & graphs & roots & roots scheduled & data scheduled & schedules executed\\
\midrule
$7$ & $7$ & $40$ & $17$ & $71$ / $242$ & $130$ (all orientations)\\
$9$ & $1{,}819$ & $10{,}118$ & $7{,}529$ & $30{,}769$ / $45{,}644$ & $76{,}474$ (all orientations)\\
\bottomrule
\end{tabular}
\end{table}

\begin{table}[ht]
\caption{Ablation at order~$9$ (roots with a schedule, of $10{,}118$).}\label{tab:ablate}
\small
\begin{tabular}{@{}lr@{}}
\toprule
moves (M), (D); end at $\le1$ token and $\le1$ live tail & $7{,}045$\\
$+$ suppression (S) (the setting of this paper; (S) is on in all rows below) & $7{,}529$\\
$+$ end as soon as $\le1$ token is left (needs Conjecture~\ref{conj:mt}) & $7{,}620$\\
$+$ end also at a sparse family of live tails (\texttt{separation\_sparse}) & $7{,}529$\\
without drops (D), end at $\le1$ token (needs Conjecture~\ref{conj:mt}) & $7{,}558$\\
without drops (D), end at $\le1$ token and $\le1$ live tail & $0$\\
\bottomrule
\end{tabular}
\end{table}

\begin{table}[ht]
\caption{Residual families of T(6)-i (\texttt{t6-dense/scripts/gen\_resid.py}: \texttt{resid\_O8}
$=$ \texttt{gen\_resid.py 6 8 2 7 SHAPE}; \texttt{resid\_O10} $=$ \texttt{gen\_resid.py 4 10 2 11
SHAPE} and \texttt{gen\_resid.py 4 8 1 77 C4}; our seeds, since those of the families
of~\cite[\S6.3]{DT} are not recorded). Suppression on. Each search: node budget $2\cdot10^6$, the
default move order (\texttt{data/runs/*\_s1\_b2e6.txt}). The roots left undecided that are not of
degree $n-2$ were run again at budget $2\cdot10^7$ with the move order shuffled by seeds $7$ and $11$
(\texttt{*\_undec\_open\_seed*.txt}; for \texttt{resid\_O10} this scheduled two more, for E28 none;
\texttt{resid\_O8} was not rerun). ``Search'': every search finished without a schedule. ``Prop.~6.1'':
undecided by the search, of degree $n-2$, so without a schedule by Proposition~\ref{prop:nminus2}. The
interim counts $71/180$ and $107/144$ reported at milestone~1 came from budget $2\cdot10^5$
(\texttt{*\_s1\_b2e5.txt}).}\label{tab:resid}
\small
\begin{tabular}{@{}llrrrrr@{}}
\toprule
 & & & & \multicolumn{2}{c}{no schedule} & \\
family & $D$ & roots & scheduled & search & Prop.~6.1 & undecided\\
\midrule
\texttt{resid\_O8}, $|O|=8$, $|J|=15$ & $\emptyset$ & $36$ & $6$ & $25$ & $2$ & $3$\\
 & $K_2$ & $36$ & $15$ & $15$ & $3$ & $3$\\
 & $P_3$ & $36$ & $19$ & $9$ & $0$ & $8$\\
 & $K_3$ & $36$ & $14$ & $10$ & $1$ & $11$\\
 & $C_4$ ($\mathrm{TW}=2$) & $36$ & $18$ & $10$ & $0$ & $8$\\
\midrule
\texttt{resid\_O10}, $|O|=10$, $|J|=17$ & $\emptyset$ & $24$ & $23$ & $0$ & $1$ & $0$\\
 & $K_2$ & $24$ & $20$ & $0$ & $4$ & $0$\\
 & $P_3$ & $24$ & $19$ & $1$ & $4$ & $0$\\
 & $K_3$ & $24$ & $21$ & $0$ & $3$ & $0$\\
 & $C_4$ ($\mathrm{TW}=2$) & $24$ & $15$ & $7$ & $2$ & $0$\\
\phantom{x}\quad and $|O|=8$, $|J|=15$ & $C_4$ ($\mathrm{TW}=1$) & $24$ & $13$ & $6$ & $2$ & $3$\\
\midrule
E28 hard graphs~\cite[E28]{EXP} & all & $48$ & $10$ & $18$ & $12$ & $8$\\
clt graph \texttt{M\textasciitilde\textasciitilde zRz\{\textasciitilde m\textasciitilde\textasciitilde lX|j\textasciitilde\_} & $\emptyset$ & $6$ & $0$ & $6$ & $0$ & $0$\\
\bottomrule
\end{tabular}
\end{table}

\section{The residual}\label{sec:resid}

\subsection*{What is left} The unscheduled roots, with the exact verdicts of Tables~\ref{tab:small}
and~\ref{tab:resid}, are listed in \texttt{data/runs/*\_none.txt} (proved: every search finished
without a schedule) and \texttt{data/runs/*\_undec.txt} (some search hit the budget).
\begin{itemize}
\item Order~$7$: $23$ of the $40$ roots. (At order $7$ with $d(z)=6$ every $J$ is a cone $K_1+B$
over an Eulerian graph $B$ of order $6$.) They are all six roots of \texttt{Fr\textasciitilde vw} ($B=K_{2,2,2}$), of
\texttt{Fqz\textasciicircum w} and of \texttt{Fut\textasciitilde w}, and the roots of degree
$5=n-2$ of \texttt{FsQ\textasciitilde w} and \texttt{Fsz\textbackslash w}. These are the smallest
unscheduled instances.
\item Order~$9$: $2{,}589$ of the $10{,}118$ roots, all proved. Every one of the $1{,}597$ roots of
degree $7=n-2$ is among them (Proposition~\ref{prop:nminus2}); the other $992$ have degree $5$ ($783$)
or $3$ ($209$). Only $260$ unscheduled roots have no legal first move in any data choice and
orientation (\texttt{firstmove.py}).
\item Residual families (Table~\ref{tab:resid}): without a schedule $69+6$ ($|O|=8$; by the search,
and by Proposition~\ref{prop:nminus2} among the roots the search left open) and $14+16$ ($|O|=10$),
undecided $33$ and $3$; E28: $18+12$ without, $8$ undecided; the clt $\emptyset$ graph: all $6$
without (search). In the
residual families a legal first move exists at $178$ of $180$, $144$ of $144$ and $48$ of $48$
roots: the obstruction is not local.
\end{itemize}

\subsection*{Roots of degree $n-2$} The degree of the root explains most of the residual.

\begin{proposition}\label{prop:nminus2}
Let $\delta(J)\ge3$, $k\ge1$ pairs, and suppose $r$ has fewer than $k$ non-neighbours in $J-r$. Then
no data choice and orientation at $r$ has a schedule. In particular, for $d(z)=6$ ($k=2$) there is
no schedule at a root of degree $|J|-2$.\lean{not\_steerK\_triangle, not\_steerK}
\end{proposition}

In Lean the hypothesis $\delta(J)\ge3$ is replaced by $d_J(w)\ge2$, which is all the proof uses.

\begin{proof}
Let $t<k$ be the number of non-neighbours of $r$ in $J-r$. For a graph $K$ on $V(M^+)$ let
\[A(K)=\{v\in V(J)-r:\ d_K(v)>0,\ v\not\sim_K r\}\]
(Lean \texttt{nonNbr}; the leaves $\ell_i$ are not in $V(J)$). Call a token move into $r$ an
\emph{entry}, and call retirements, suppressions and drops that isolate the stem \emph{removals}.

\emph{Before the first entry.} In $M^+$ the vertex $r$ has all its edges of $J$ (no deleted edge
of $M$ is at $r$), so $|A(M^+)|\le t$; $r$ is odd and is not a token, a leaf or a stem, and $N(z)=\{r\}$.
Until an entry the degree of $r$ does not change: a move out of $r$ would need a token at $r$, a drop
is at a stem $y_i\ne r$, and a suppression keeps the degrees of $a,b$. So $r$ stays odd and is not a
token, and the only edge at $r$ that can be deleted is $pr$ in a suppression of $p$ with
$a=r$, after which $p$ is isolated. Also, an isolated vertex stays isolated (the only edge ever added,
$ab$ in (S), joins two neighbours of $p$). Hence no move adds a vertex to $A(K)$, and $N(z)=\{r\}$
throughout. We claim that every removal takes a vertex out of $A(K)$.
\begin{itemize}
\item A retirement $p\to u$: before it $N(u)=\{p\}$ with $p\ne r$, so $u\not\sim r$; $u$ is not a
leaf (live leaves are in $X$, dead leaves are isolated) and $u\ne z$, so $u\in A(K)$; afterwards
$u$ is isolated.
\item A suppression of $p$ with $N(p)=\{a,b\}$; split on the \emph{current} adjacency of $p$ to $r$.
If $p\not\sim r$, then $p\in A(K)$ ($p$ is a token, so $p\ne r,z$ and $p$ is not a leaf), and
afterwards $p$ is isolated. If $p\sim r$, then $r\in\{a,b\}$, say $a=r$; then $b\not\sim r$ (as
$a\not\sim b$), $b$ is not a leaf, $b\ne z$, so $b\in A(K)$, and afterwards $b\sim r$.
\item A drop that isolates its stem $y_i$: then $N(y_i)=\{\ell_i\}$, so $y_i\in A(K)$, and
afterwards $y_i$ is isolated.
\end{itemize}
So after $R$ removals, $|A(K)|\le t-R$. The number of tokens is at least $k+1-R$: it starts at $k+1$
($w$ and the $y_i$; $d_M(w)=d_J(w)-1\ge2$), a non-retiring move keeps it, a drop does not lower it,
and only retirements and suppressions lower it, by one each.

\emph{There is no entry.} At the first entry $p\to r$, every other token $q$ is not adjacent to $r$
after the deletion of $pr$, hence not before, and $q\in A(K)$. So $k+1-R-1\le t-R$, i.e.\ $k\le t$,
which is false.

So the whole schedule lies before the first entry and has at most $t<k\le2k-1$ removals, while
Proposition~\ref{prop:potential} requires $2k-1$ ($\Phi(M^+)=2k+1$, as $d_J(w)\ge3$). No schedule
exists. (Lean \texttt{not\_steerK} carries the inequalities $|A(K)|+2\le|\tau(K)|$ and
$|A(K)|+3\le\Phi(K)$ through every move, with the same case split in \texttt{obstr\_supp}.)
\end{proof}

The data agree: at every root of degree $n-2$ that the search proves unscheduled, the least
potential reached is $5$ or $4$, never less (\texttt{explore.py}, exhaustive over all reachable states
of every data choice and orientation; its state cap was hit at three roots of degree $n-2$ of
\texttt{resid\_O8}, where the uncapped exploration of audit A36 gives the same values). For $d(z)\ge6$ such roots therefore need a different method; for $d(z)=8$ the
proposition excludes roots with at most two non-neighbours.

\subsection*{The rest: frozen and partial roots} At every root that the search proves unscheduled
(the roots settled only by Proposition~\ref{prop:nminus2} are not explored), \texttt{explore.py}
explores all reachable states of every data choice and orientation and records the least potential
$\Phi$ reached (Proposition~\ref{prop:potential}; $\Phi=5$ at the start, a schedule ends at
$\Phi\le2$ with at most one token and one live tail). Call a root \emph{frozen} if $\Phi$ never drops,
so that no removal is ever possible.
\begin{center}\small
\begin{tabular}{@{}lrrrrrr@{}}
\toprule
 & unscheduled & degree $n-2$ & other frozen & other, least $\Phi=4$ & $=3$ & $=2$\\
\midrule
order $7$ & $23$ & $20$ & $0$ & $1$ & $0$ & $2$\\
order $9$ & $2{,}589$ & $1{,}597$ & $113$ & $597$ & $121$ & $161$\\
\texttt{resid\_O8} & $69$ & $36$ & $27$ & $4$ & $2$ & $0$\\
\texttt{resid\_O10} & $14$ & $2$ & $7$ & $3$ & $2$ & $0$\\
E28 hard & $18$ & $12$ & $6$ & $0$ & $0$ & $0$\\
clt graph & $6$ & $4$ & $2$ & $0$ & $0$ & $0$\\
\bottomrule
\end{tabular}
\end{center}
(At order~$7$, $15$ of the $20$ roots of degree $n-2$ are frozen; at order~$9$, $376$ of $1{,}597$.)
A least potential of $2$ without a schedule means two tokens and no live tail, or no token and two
live tails; the second kind is what Conjecture~\ref{conj:mt} would settle. In the dense graphs most
unscheduled roots are frozen: of the $53$ search-proved unscheduled roots of degree less than $n-2$ in
the residual families, the E28 graphs and the clt graph, $42$ are frozen. Of these $42$, $38$ have
exactly three non-neighbours of $r$ and four have five, so Proposition~\ref{prop:nminus2} does not
reach them. Like Proposition~6.10 of~\cite{TRN} for one pair, they show that a
dense $N(r)$ blocks steering: the root is the only way to the leaf $z$, and it cannot be entered while
other tokens are near it.


\begin{conjecture}[multi-tail, $\MT$]\label{conj:mt}
If $K$ has at most one even vertex and $T$ satisfies (TS) and (TI), then $K$ is good, whatever the
number of live tails.
\end{conjecture}

$\MT$ holds in all our tests (\texttt{scripts/mt\_conj.py}; the oracle \texttt{route\_ext}, query
\texttt{S}). An instance is a connected graph $G$ with leaves attached at distinct stems, and avoided
vertices that are distinct, odd, outside the stems and non-adjacent to their own stem, such that the
result has at most one even vertex; so the test covers only distinct stems and distinct avoided
vertices, while (TS) allows both to be shared. With two tails, every connected $G$ of order $n\le9$
($1$, $29$, $113$, $6{,}502$, $18{,}858$ and $2{,}860{,}182$ unordered instances for $n=4,\dots,9$);
with three tails, every connected $G$ of order $n\le7$ ($92$ unordered instances at $n=6$ and
$1{,}026$ at $n=7$; \texttt{mt\_conj.py} counts ordered triples, $552$ and $6{,}156$). Audit A36 extended
the test with its own solver to everything (TS) allows, shared stems and shared avoided vertices
included: no failure among $37{,}707$ two-tail instances ($n\le8$) and $17{,}213$ three-tail
instances ($n\le7$). If $\MT$ is true, the
end of a schedule can be taken at $\le1$ token; at order~$9$ this adds $91$ roots
(Table~\ref{tab:ablate}).

\begin{remark}[where an induction for $\MT$ stops]\label{rem:mtind}
The leaf-tail lemma is $\MT$ for one tail. Run its induction with two live tails. With no even
vertex, drop a tail and apply the leaf-tail lemma. If the even vertex $g$ has a neighbour outside
$X$, take Lov\'asz's Case~I step there (Lemma~\ref{lem:fan} with one token). Otherwise every
neighbour of $g$ lies in $X$. Assume distinct stems (if $s_1=s_2=g$ and $N(g)=\{c_1,c_2\}$, the
component $c_1gc_2$ is one path, which avoids $x_1,x_2$, and the claim is immediate). A stem is
adjacent to no leaf but its own, and $x_i\not\sim s_i$, so
$N(g)=\{x_1,x_2\}$ or $N(g)=\{c_j,x_i\}$ with $g=s_j$, $i\ne j$. If $N(g)=\{x_1,x_2\}$ with $x_1\not\sim x_2$,
suppress $g$ and drop. The two configurations left, $g$ on a triangle $gx_1x_2$ and a stem
$s_j$ of degree two with neighbours $c_j,x_i$, are where we do not know how to proceed.
\end{remark}

\section{The all-odd lift}\label{sec:aodd}

The coordinator observed a second lift to Test1~\cite[\S31]{SYN}, whose base lives in a graph with
no even vertex at all. We check it, run steering on it, and find that steering does worse there.

\subsection*{The lift} Let $z$ be the only even vertex of $J$, $d(z)=2k$, $r\in N(z)$, and let
$\Phi$ be a perfect matching of $J[N(z)]$ by edges of $J$, with $ra_0\in\Phi$. Put $B=J-z$ and
$H=B-\Phi$.

\begin{proposition}[the all-odd lift~{\cite[\S31]{SYN}}]\label{prop:aoddlift}
$H$ has no even vertex. If a Lov\'asz decomposition of $H$ has $r\notin V(Q(a_0))$ and, for every
other pair $ab\in\Phi$, $a\notin V(Q(b))$ or $b\notin V(Q(a))$, then Test1 holds at $(J,z,r)$.
\end{proposition}

\begin{proof}
A vertex of $N(z)$ loses its edge to $z$ and one edge of $\Phi$; every other vertex keeps its
$J$-degree, which is odd. So every vertex of positive degree in $H$ is odd, every vertex is an end of
exactly one path, and there are $(|J|-1)/2$ paths. Orient each pair $(y,x)$ so that $x\notin
V(Q(y))$, with $(y,x)=(a_0,r)$, and append $x$ to $Q(y)$ at $y$. Distinct pairs have disjoint
vertices, so each path receives at most one vertex at each of its ends, and the vertex appended at
one end is not on the path or appended at the other end. The result decomposes $B$ into $(|J|-1)/2$
paths, and $r$ is an end of $Q(r)$ and of $Q(a_0)+r$, which are different paths since $r\notin
Q(a_0)$.
\end{proof}

The executor (\texttt{mt\_verify.py}, lines \texttt{SCHEDA}) performs exactly this append and checks
the Test1 decomposition. The lift is not formalised.

\subsection*{Steering on the all-odd base} Hang a leaf $\ell_0$ at $a_0$ with avoided vertex $r$ (a
fixed orientation), and a leaf $\ell_i$ at one vertex $y_i$ of each other pair, with avoided vertex
$x_i$ its partner. The tails satisfy (TS) and (TI), and the tokens are exactly the $k$ stems:
there is no $w$, and no leaf $z$ (now isolated). Theorem~\ref{thm:transfer} applies verbatim: a
schedule gives a Lov\'asz decomposition of $H$ with $x_i\notin V(Q(y_i))$ for every $i$, hence Test1
by Proposition~\ref{prop:aoddlift}. The potential starts at $\Phi=2k$, so a schedule needs $2k-2$
removals (four at $d(z)=6$, against three for the triangle base).

\begin{proposition}[the obstruction, all-odd form]\label{prop:aoddobs}
Let $\delta(J)\ge3$, $d(z)=2k\ge4$, and suppose $r$ has at most $k-2$ non-neighbours in $J-r$.
Then no perfect matching $\Phi$ and no orientation give a schedule. At $d(z)=6$ this is every root
of degree $|J|-2$.\lean{not\_steerK (abstract form)}
\end{proposition}

\begin{proof}
Apply the invariant of \texttt{not\_steerK}. Let $A=A(K)$ be the current set of non-neighbours of
$r$ of positive degree that are not leaves (as in the proof of Proposition~\ref{prop:nminus2}), $\rho$
the number of live tails whose avoided vertex is $r$, and
$\Phi=\#\text{tokens}+\#\text{live tails}$; the claim is that
\[|A|+2\le\#\text{tokens}+\rho\qquad\text{and}\qquad|A|+3\le\Phi\]
are kept by every move and contradict the end ($\Phi\le2$). At the start $r$ is odd ($d_H(r)=d_J(r)-2$),
neither a leaf nor a stem; $A$ consists of $a_0$ (no longer adjacent to $r$) and the $t$
non-neighbours of $r$ in $J$ ($z$ is isolated, and nothing else at $r$ is deleted), so $|A|=t+1$;
there are $k$ tokens, $k$ live tails and $\rho=1$. So the inequalities read $t+3\le k+1$ and $t+4\le2k$,
which hold for $t\le k-2$. Invariance: a token cannot enter $r$, since while $\rho>0$ the vertex $r$ is
in $X$, and when $\rho=0$ every other token would be a non-neighbour of $r$, i.e.\ in $A$, against the
first inequality. Every other move is as in the proof of Proposition~\ref{prop:nminus2} (with the
suppression case split on the current adjacency of $p$ to $r$): no move adds a vertex to $A$, a
removal takes one out and lowers $\#\text{tokens}+\rho$ and $\Phi$ by at most one; the drop of the
tail $(\ell_0,a_0,r)$ lowers $\rho$ by one but adds a token or uses up $a_0$. (Lean
\texttt{not\_steerK} proves exactly this invariance, for any graph, tails and $r$.)
\end{proof}

\begin{table}[ht]
\caption{Steering on the all-odd lift against the triangle base (roots with a schedule / roots with
data). Orders $7$ and $9$: exact. Residual families, all-odd column: node budget $2\cdot10^6$, default move order
(\texttt{*\_aodd\_b2e6.txt}); triangle-base column as in Table~\ref{tab:resid}.}\label{tab:aodd}
\small
\begin{tabular}{@{}lrr@{}}
\toprule
 & triangle base & all-odd lift\\
\midrule
$d(z)=4$, order $7$ & $73$ / $74$ & $23$ / $52$\\
$d(z)=4$, order $9$ & $11{,}347$ / $11{,}389$ & $6{,}456$ / $7{,}572$\\
$d(z)=6$, order $7$ & $17$ / $40$ & $0$ / $36$\\
$d(z)=6$, order $9$ & $7{,}529$ / $10{,}118$ & $2{,}101$ / $8{,}340$\\
\quad with the end of Conjecture~\ref{conj:mt} & $7{,}620$ & $2{,}725$\\
\texttt{resid\_O8} & $72$ / $180$ & $14$ / $180$\\
\texttt{resid\_O10} & $111$ / $144$ & $59$ / $144$\\
E28 hard graphs & $10$ / $48$ & $0$ / $48$\\
\bottomrule
\end{tabular}
\end{table}

At order $9$ no root of degree $n-2$ has an all-odd schedule ($0$ of $1{,}546$), as
Proposition~\ref{prop:aoddobs} predicts, and no root at which no vertex of $N(z)$ has degree three ($0$
of $2{,}430$; computed, not proved). All $5{,}700$ all-odd schedules at order~$9$ (every data choice and
orientation) were executed by \texttt{mt\_verify.py}. The coordinator's exhaustive search finds an
all-odd base at every root tested~\cite[\S31]{SYN}, so the bases exist; steering does not reach
them. The reason is the count: with one tail more and no $w$, one more removal is needed (four
instead of three at $d(z)=6$), and the leaf $z$, the one retirement spot the triangle base always
has, is gone. In the residual families the all-odd search proves that no schedule exists at $136$
of the $180$ roots of \texttt{resid\_O8} (all $36$ roots of its $C_4$ shape), $21$ of the $144$ of
\texttt{resid\_O10}, and $28$ of the $48$ roots of the E28 graphs; every all-odd schedule found there was executed by
\texttt{mt\_verify.py}. A coverage theorem in the all-odd setting is not in sight with this move
set.

\section{Formalisation}\label{sec:lean}

The files are \texttt{lean/GallaiLean/MultiToken.lean} and \texttt{MultiTokenObstruct.lean}, with no
\texttt{sorry}. \texttt{lake build} is green ($18{,}428$ jobs), and \texttt{\#print axioms} gives only
\texttt{propext}, \texttt{Classical.choice} and \texttt{Quot.sound} for every theorem and lemma in the table below
($22$ in all; \texttt{scripts/axioms\_mt.lean}, output \texttt{data/axioms.txt}).

\begin{center}\small
\begin{tabular}{@{}ll@{}}
\toprule
manuscript & Lean\\
\midrule
Definition~\ref{def:tails} & \texttt{TailSys}, \texttt{TailInv}, \texttt{TailGood},
\texttt{OffTails}, \texttt{OneLive}, \texttt{SuppOK}\\
Definition~\ref{def:sched} & \texttt{SteerK} (constructors \texttt{fin}, \texttt{move},
\texttt{drop}, \texttt{supp})\\
Lemma~\ref{lem:fan} & \texttt{fan\_step\_multi}\\
Lemma~\ref{lem:inv} & \texttt{evenSet\_move/drop/supp}, \texttt{tailInv\_move/drop/supp},
\texttt{live\_move/drop/supp}\\
Lemma~\ref{lem:steps} & \texttt{tailGood\_move}, \texttt{tailGood\_drop}, \texttt{tailGood\_supp},
\texttt{splitOff\_supp}\\
Lemma~\ref{lem:mtl} & \texttt{tailGood\_fin}\\
Theorem~\ref{thm:transfer} & \texttt{steerK\_good}\\
Theorem~\ref{thm:tb} & \texttt{triangleBase\_of\_steerK} (the ``in particular''; with
\texttt{triTails\_sys}, \texttt{triTails\_inv})\\
Corollary~\ref{cor:test1} & \texttt{test1At\_of\_steerK}\\
Proposition~\ref{prop:nminus2} & \texttt{not\_steerK\_triangle}; the invariant \texttt{ObstrInv},
\texttt{not\_steerK}\\
Proposition~\ref{prop:aoddobs} & \texttt{not\_steerK} (abstract form; the instance on paper)\\
\bottomrule
\end{tabular}
\end{center}

Faithfulness. \texttt{SteerK} quantifies over an arbitrary token set $\tau$; the theorems use it
with $\tau$ the even set of the graph, which is what \texttt{mtok.cpp} and
\texttt{mt\_verify.py} compute. The C++ move conditions are exactly the constructor hypotheses:
(M) $u$ odd ($\Leftrightarrow u\notin\tau$, as $u\sim p$), $u\notin X$, no token neighbour of $u$
other than $p$ after the deletion; (D) a live tail with odd stem; (S) degree two, non-adjacent
neighbours, not the leaf of any tail (a dead leaf is isolated, so this is the C++ test on live
leaves), and not the pair $x_is_i$ of a live tail; the end: $\le1$ even vertex and $\le1$ live
tail. The Lean form of Theorem~\ref{thm:tb} exports its last clause, \texttt{TriangleBase J z r};
the first clause ($x_i\notin V(Q(y_i))$ for every $i$) is proved inside the proof, as the disjunct
\texttt{Or.inl} for every pair, but not stated. Audit A36 ran Example~\ref{ex:short} through the Lean
pipeline to \texttt{Test1At} by \texttt{decide} (\texttt{research/audit/A36-scripts/ex33\_steerK.lean},
which adds the \texttt{Decidable} instances needed for \texttt{OffTails}, \texttt{SuppOK},
\texttt{OneLive} and nested graph operations). Theorem~\ref{thm:cover},
Proposition~\ref{prop:aoddlift} and the all-odd instance of Proposition~\ref{prop:aoddobs} are not
formalised.

\section{What remains}\label{sec:open}

\begin{enumerate}
\item The dense residual. With this move set, schedules exist at $183$ of the $324$ residual-family
roots; at $|O|=8$ every shape has roots without any schedule (Remark~\ref{rem:noshape}), and no root
of degree $n-2$ has one (Proposition~\ref{prop:nminus2}). A proof of the triangle base there needs a
new token-removing move, or an end lemma with two even vertices. At $|O|=10$, in the shapes
$\emptyset$, $K_2$ and $K_3$ of our sample, every root has a schedule or has degree $n-2$
(Table~\ref{tab:resid}); whether this holds for these shapes in general is open.
\item Conjecture~$\MT$; it is the multi-tail generalisation of the leaf-tail lemma.
\item The all-odd lift (Section~\ref{sec:aodd}): its bases exist in all tests, but steering covers a
quarter of the roots at order $9$; all-odd DL$(k)$~\cite[\S31]{SYN} needs another tool.
\item The frozen roots of degree less than $n-2$ (Section~\ref{sec:resid}) are computed; a proof of the
obstruction there, extending Proposition~\ref{prop:nminus2}, is open.
\end{enumerate}

\begin{thebibliography}{99}
\bibitem{EXP} The experiments log, \texttt{research/experiments.md}, entry E28 (E): the eight hard
graphs of the dense residual.
\bibitem{DT} The direct lift to Test1 through triangles at $z$, \texttt{manuscripts/direct-test1}.
\bibitem{Fan05} G.~Fan, Path decompositions and Gallai's conjecture, \emph{J. Combin. Theory Ser.
B} \textbf{93} (2005), 117--125.
\bibitem{L68} L.~Lov\'asz, On covering of graphs, in: \emph{Theory of Graphs} (Proc. Colloq.,
Tihany, 1966), Academic Press, 1968, 231--236.
\bibitem{NA} Neighbour avoidance, \texttt{manuscripts/na}.
\bibitem{RD} Route descent, \texttt{manuscripts/route-descent}.
\bibitem{SYN} The project synthesis, \texttt{research/synthesis.md}: \S28 (state of the programme), \S31
(the all-odd lift, coordinator).
\bibitem{T6} The dense T(6)-i residual, \texttt{manuscripts/t6-dense}.
\bibitem{TRN} The tight transfer TRN, two-token steering, \texttt{manuscripts/trn}.
\end{thebibliography}
\end{document}
